% ============================================
% LUC-09
% Gerado de molambudos.md — build_miolo.py
% ============================================
\fragtitle{luc}{L}{LUC}{09}{A Exumação de Oliveira}
\markboth{LUC-09 --- A Exumação de Oliveira}{}

\noindent \textit{47 anos depois, o crânio na vala}

\noindent \textbf{Junho de 2026 --- Quixeramobim, Ceará}

\vspace{0.5em}\noindent\rule{\textwidth}{0.5pt}\vspace{0.5em}

\noindent Lúcia foi ao Ceará no dia 13 de junho de 2026. Exatos 47 anos depois da morte de Joaquim e do desaparecimento de Oliveira.

\noindent Ela não sabia o que esperava encontrar. A vala estava lá --- a clareira de terra preta, a árvore torta, o cheiro adocicado que subia do chão como se a terra estivesse respirando.

\noindent Ela se ajoelhou.

\noindent Enterrou as mãos na terra preta.

\noindent A terra estava quente.

\noindent \textit{\textquotedblleft{}Dr. Oliveira?\textquotedblright{}\textit{, ela chamou. }\textquotedblleft{}O senhor está aqui?\textquotedblright{}}

\noindent O vento respondeu.

\noindent \textit{\textquotedblleft{}Eu li o seu diário. Li o de Joaquim. Eu sei o que aconteceu. Eu estou sentindo o mesmo. O olho. O cheiro. Os sonhos. Está acontecendo comigo também.\textquotedblright{}}

\noindent Ela tirou uma pá da mochila.

\noindent Suas mãos tremiam, mas ela não parou. Tinha medo do que ia encontrar --- e mais medo ainda de não encontrar nada, de descobrir que também estava enlouquecendo sozinha, como Joaquim, como Oliveira antes dela.

\noindent E começou a cavar.

\vspace{0.5em}\noindent\rule{\textwidth}{0.5pt}\vspace{0.5em}

\noindent Cavou por duas horas.

\noindent A terra era fofa, oleosa. Fácil de cavar. Como se já tivesse sido remexida muitas vezes.

\noindent A um metro de profundidade, a pá bateu em algo duro.

\noindent Ela cavou com as mãos.

\noindent Era um osso.

\noindent Não de animal. Humano.

\noindent Ela cavou mais. Uma costela. Uma vértebra. Um crânio.

\noindent O crânio tinha uma fissura na têmpora direita. Como de um objeto pontiagudo.

\noindent E dentro da órbita direita --- a órbita ainda estava escura, ainda guardava resquícios de tecido --- algo brilhou.
\noindent Um reflexo amarelo, fraco, no fundo da órbita.
\noindent Podia ser a lanterna. Podia ser mica no barro.

\noindent Ela aproximou o rosto.

\noindent Um reflexo amarelo.

\noindent O olho ainda brilhava.

\noindent Mesmo depois de 47 anos.

\noindent Lúcia sentou no chão da vala, com o crânio de Oliveira na mão, o gosto de açúcar queimado nos pulmões, e uma certeza subindo pelo tornozelo.

\noindent Era uma corrente de corpos.

\noindent E ela era a próxima.

\vspace{0.5em}
\noindent\textit{\textbf{↪ Rotas:} \rota{LUC-10} \rotasep \rota{DOC-15}}